% Source: Wald GR, physical PDF pp. 22--23
%         (printed pp. 9--10).
% Coverage: Chapter 1 problem 1, Car and garage paradox.
% Figures/tables/footnotes: no figures, tables, or footnotes.
% Status: source-reviewed and compile-clean.
% Uncertainties: none.

\subsection*{Problem}
\addcontentsline{toc}{subsection}{Problem}

\noindent\textbf{1.}\label{prob:1.1}\quad\emph{Car and garage paradox:} The lack of
a notion of absolute simultaneity in special relativity leads to many supposed
paradoxes. One of the most famous of these involves a car and a garage of equal
proper length. The driver speeds toward the garage, and a doorman at the garage is
instructed to slam the door shut as soon as the back end of the car enters the garage.
According to the doorman, ``the car Lorentz contracted and easily fitted into the garage
when I slammed the door.'' According to the driver, ``the garage Lorentz contracted and
was too small for the car when I entered the garage.'' Draw a spacetime diagram
showing the above events and explain what really happens. Is the doorman's statement
correct? Is the driver's statement correct? For definiteness, assume that the car crashes
through the back wall of the garage without stopping or slowing down.
